\scanpage{41}
\mspage{15}

Si $X$ est une surface conforme hyperbolique, alors la métrique canonique hyperbolique sur $X^{\mathrm{or}}$ se descend en une métrique hyperbolique sur $X$, pour laquelle le bord est géodésique.

Plus généralement, chaque fois que $X^{\mathrm{or}}$ est munie d'une métrique canonique, invariante par autom conformes [ou seulem$^{\mathrm t}$ par antiinvolutions, ce qui est le cas aussi si $X^{\mathrm{or}}\simeq\mathbb C$] alors celle-ci se descend à $X$, avec bord géodésique.

Quand $X$ est \deleted{\unsure{avec}}sans bord, \deleted{et \unsure{orientable|si elle}} alors \deleted{\unsure{la}} la métrique obtenue sur $X$ est celle déjà envisagée -- \deleted{\unsure{est}}\inserted{il y en a} dans tous les cas, sauf $X\simeq\mathbb P^1_{\mathbb C}$. On fera attention cependant que lorsque $X$ est le plan projectif réel, la deuxième prescription \deleted{\unsure{sans}} plus ne fournit\deleted{\unsure{s}} \uline{pas} de métrique -- elle le fait si on exige que la métrique sur $X^{\mathrm{or}}$ soit en $R(X^{\mathrm{or}})$, \uline{et} stable par $\sigma$ (ce qui la détermine sans ambiguïté). \deleted{De =}\inserted{\deleted{\unsure{De faç|De fait}}En sens inverse}, dans le cas $X\simeq\mathbb C$, $X^{\mathrm{or}}\simeq\mathbb C\amalg\mathbb C$, le procédé direct donne une classe de métriques (mod cte), la deuxième
% archive sheet number: 43
